\documentclass[10pt,a4paper]{amsart}
\usepackage[margin=19mm]{geometry}
\usepackage{amsmath,amssymb,mathtools}
\usepackage[scr=rsfs,cal=euler]{mathalfa}
\usepackage{xfrac}
\usepackage[unicode,hidelinks,bookmarks=false]{hyperref}
\newcommand{\ie}{i.e.\ }
\newcommand{\cf}{cf.\ }
\newcommand{\resp}{respectively\ }
\newcommand{\Subsection}{\subsection}
\providecommand{\nobreakdash}{\nobreak\hbox{-}}
\setlength{\emergencystretch}{2em}
\multlinegap0pt
\usepackage{amssymb}
\usepackage{cancel}
\usepackage{xcolor}
\newtheorem{theorem}{Theorem}
\def \LL {\mathcal{L}}
\def \R {\mathbb{R}}
\def \RN {\mathbb{R}^N}
\def \e {\varepsilon}
\title{Blow-up of solutions to semilinear parabolic equations driven by mixed local-nonlocal operators with large initial data}
\author{Stefano Biagi, Fabio Punzo, Eugenio Vecchi}
\date{Source: arXiv:2605.07911v1 (CC BY 4.0)}
\begin{document}
\maketitle

\noindent\emph{Source and licence.} Blow-up of solutions to semilinear parabolic equations driven by mixed local-nonlocal operators with large initial data. Version arXiv:2605.07911v1 (CC BY 4.0), distributed by arXiv under the Creative Commons Attribution 4.0 International licence, http://creativecommons.org/licenses/by/4.0/, as recorded in the arXiv licence metadata for this exact version and shown on the abstract page https://arxiv.org/abs/2605.07911v1. Full licence notice: \texttt{licenses/arxiv-2605.07911-CC-BY-4.0.txt}.\par\smallskip

\noindent\textit{Editorial scope.} Counted unit: the article's theorem teo3 (source0.tex lines 956--960). The article's complete original proof of it is delivered with it (source0.tex lines 961--979, one \texttt{proof} environment of 19 lines). No mathematics is written, repaired or re-derived: the statement and the proof are the article's own lines, in the article's own notation, and no retained line is substituted or re-flowed.  Retained uncounted as the source-local context the proof needs: lines 136--141; lines 587--589; lines 886--888; lines 929--937.  Nothing was omitted from any retained span, so the package carries no omission record.  No retained line contains a citation, so the wrapper carries no bibliography and no bibliography entry.  No retained line contains a figure environment, an image-inclusion command or drawing code, so no image is delivered and the package declares no asset dependency.  Editorial formatting only, in addition: an AMS \texttt{amsart} preamble that restates the article's own declarations and macros used by the retained lines, namely the environments \texttt{theorem} on separate counters, together with the standard packages those lines need.  The retained environments print in this document as Theorem 1 in order of appearance, whereas the article prints the same items with the numbering of its own full document.  The complete native source of the article as distributed in the arXiv e-print for this exact version is bundled unchanged as \texttt{sources/200120/source0.tex}.
\begin{equation}\tag{$\mathrm{CP}_{\LL}$}\label{eq:mainPbParabolicIntro}
	\begin{cases}
	\partial_t u+\mathcal{L}	u = f(u) & \text{in $\mathbb{R}^N\times(0,+\infty)$} \\
	u(\cdot,0) = u_0 & \text{in $\mathbb{R}^N$},
	\end{cases}
\end{equation}

\begin{equation} \label{eq:sflambdapower}
\mathbf{s}_f(\lambda) = \lambda^{1/(p-1)}.
\end{equation}

\begin{equation} \label{eq:deflambdazero}
\lambda_0 =  \max\{\lambda_1,\lambda_2\}> 0.	
\end{equation}

\begin{theorem} \label{thm:NonexistenceExplicit}
	Let $u_0\in C_+(\mathbb{R}^N)\cap L^\infty(\mathbb{R}^N)$, and assume that there exist $\beta > N/2,\,\e\in(0,1]$ such that the following
	estimate holds
	\begin{equation} \label{eq:conditionKaplanExplicit}
		\frac{\e^{N/2}}{\mathbf{c}_\beta}\int_{\R^N}\frac{u_0(x)}{(1+\e|x|^2)^\beta}\,dx >\mathbf{s}_f(\e^s\lambda_0)
	\end{equation}
	\emph{(}where $\lambda_0 > 0$ as in \eqref{eq:deflambdazero}, and
	$\mathbf{c}_\beta = \|\Psi_\beta\|_{L^1(\R^N)}$\emph{)}. Then, the Cauchy problem \eqref{eq:mainPbParabolicIntro}, with initial datum $u_0$, does not admit global classical solution.
\end{theorem}

\begin{theorem}\label{teo3}
Let $f(z) = z^p$, with $1<p<p_F$. Then, the Cauchy problem \eqref{eq:mainPbParabolicIntro}
\emph{(}associated with this $f$\emph{)} does not admit global classical solutions \emph{for every non-zero initial datum}
$$u_0\in C_+(\mathbb{R}^N)\cap L^\infty(\mathbb{R}^N).$$
\end{theorem}

\begin{proof}
	Let $u_0\in C_+(\mathbb{R}^N)\cap L^\infty(\mathbb{R}^N),\,u_0\not\equiv 0$. By Theorem \ref{thm:NonexistenceExplicit}, it suffices to prove that there exist $\beta > N/2$ and $\varepsilon\in(0,1]$ such that\begin{equation}\label{eq:toProveSubcritical}
\begin{split}
& \frac{\e^{N/2}}{\mathbf{c}_\beta}\int_{\R^N}\frac{u_0(x)}{(1+\e|x|^2)^\beta}\,dx >\mathbf{s}_f(\e^s \lambda_0)\\
&\qquad\qquad\Longleftrightarrow\,\,
	\int_{\R^N}\frac{u_0(x)}{(1+\e|x|^2)^\beta}\,dx >\mathbf{c}_\beta \lambda_0^{1/(p-1)}\e^{s/(p-1)-N/2},
\end{split}	
\end{equation}
where we have used  that $\mathbf{s}_f(\lambda) = \lambda^{1/(p-1)}$ when $f(z) = z^p$ (see \eqref{eq:sflambdapower}).
\vspace{0.1cm}

Fix now $\beta > N/2$. On the one hand, since $u_0\geq 0$ and $u_0\not\equiv 0$, by the Monotone Convergence Theorem we have the following computation
$$\lim_{\e\to 0^+}\int_{\R^N}\frac{u_0(x)}{(1+\e|x|^2)^\beta}\,dx =
\int_{\RN}u_0(x)\,dx > 0;$$
on the other hand, since $1<p<p_F$, we have
$$\lim_{\e\to 0^+}\e^{s/(p-1)-N/2} = 0.$$
As a consequence, since both $\mathbf{c}_\beta$ and $\lambda_0$ are  independent of $\e$, there exists $0<\e^*<1$ such that
condition \eqref{eq:toProveSubcritical} is satisfied. This ends the proof.
\end{proof}

\end{document}
